\renewcommand*{\authorOfNextLines}{Helmut Quinz}
\begin{quote}
	Und jedem Anfang wohnt ein Zauber inne,\\
	Der uns beschützt und der uns hilft, zu leben.\\
	\small{Hermann Hesse, Stufen}
\end{quote}
Wurde eine Projektidee entwickelt und ihre Relevanz positiv beurteilt gilt es das Projekt zu starten.\\
Dazu ist aus Projektidee zunächst das Projektziel und der Projektumfang \textbf{exakt} zu definieren.
Beides soll \textbf{SMART} erfolgen:
\begin {table} [ht]
\stretchTableRows{1.2}
\begin{tabular}{p{5mm}p{22mm}p{20mm}p{97mm}}
		& \textbf{Englisch}	& \textbf{Deutsch} 	& \textbf{Bedeutung}\\
	S	& Specific			& Spezifisch		& Das Thema muss so präzise wie möglich definiert sein. Vage Aussagen sind zu vermeiden\\
	M	& Measurable		& Messbar			& Kriterien zur Bewertung der Erfüllung oder des Erfüllungsgrads einer Vorgabe sind festzulegen\\
	A	& Accepted			& Akzeptiert, Attraktiv
												& Das Themas sollte von allen Beteiligten und all jenen die Einfluss auf das Projekt nehmen können als sinnvoll betrachtet werden\\
	R	& Realistic			& Realistisch		& Das Thema sollte mit den zur Verfügung stehenden Ressourcen unter einer ausreichenden Wahrscheinlichkeit umsetzbar sein\\
	T	& Time-bound		& Terminiert		& Der späteste Termin zur Erfüllung der Vorgabe ist festzulegen.
\end{tabular}
\unstretchTableRows
\caption {Beschreibung der Methode \textbf{SMART}}
\end {table}
\FloatBarrier
Aus dem Projektziel(en) und dem erkannten Umfang lässt sich in Folge ein Projektauftrag formulieren.
Für eine erste grobe Planung des Projektes sind bereits hier die Meilensteine zu definieren.
Das sind bedeutende Ereignisse eines Projekts einen wichtigen Fortschritt im Projekt markierend.
Sie dienen dazu, den Projektfortschritt zu überwachen, den Zeitplan zu strukturieren und den Erfolg des Projekts zu bewerten. 
Beispiele für Meilensteine könnten sein:
\begin{itemize}
	\item Abschluss einer detaillierten Projektplanung.
	\item Fertigstellung von Fertigungsplänen
	\item Beginn der Montagearbeiten
	\item Inbetriebnahme
	\item Beginn des Probebetriebs
	\item Fertigstellung des Produkts
	\item Abschluss des Projekts
\end{itemize}
\paragraph{Wichtig} Unter anderem sind die hier gewonnenen Erkenntnisse in der DA-Datenbank bei Anlegen der Diplomarbeit anzugeben.
(\SeeRefNamePage{par:KurzfassungDiplomarbeitErlaeuterung})